% ========================== MODELO DE NEO4J ==========================
\section{Modelo de Grafo en Neo4j}
\label{sec:neo4j}

Neo4j es la capa de \textbf{relaciones}. El grafo representa quién interactúa con qué repositorio y a qué organización pertenece cada repositorio. Las restricciones se crean con \repo{scripts/init_neo4j.cypher}{scripts/init\_neo4j.cypher}.

\subsection{Nodos y Relaciones}

\begin{figure}[H]
    \centering
    \begin{tikzpicture}[
        font=\small,
        nodo/.style={circle, draw=black, thick, minimum size=2.1cm, align=center},
        arrow/.style={-{Stealth[scale=1.2]}, thick},
        lbl/.style={font=\scriptsize, align=center, fill=white, inner sep=2pt}
    ]
    \node[nodo, fill=blue!14] (u) {\textbf{:User}\\ \texttt{login}};
    \node[nodo, fill=green!16, right=5.6cm of u] (r) {\textbf{:Repo}\\ \texttt{name}\\ \texttt{owner}};
    \node[nodo, fill=orange!20, right=3.6cm of r] (o) {\textbf{:Org}\\ \texttt{login}};
    \draw[arrow] (u) -- node[lbl] {\texttt{PUSHED} $|$ \texttt{OPENED\_PR} $|$ \texttt{OPENED\_ISSUE}\\ \texttt{STARRED} $|$ \texttt{FORKED}\\[0.2em] \{\texttt{count}, \texttt{first\_seen}, \texttt{last\_seen}\}} (r);
    \draw[arrow] (r) -- node[lbl] {\texttt{BELONGS\_TO}} (o);
    \end{tikzpicture}
    \caption{Modelo de grafo de propiedades: tres etiquetas de nodo y seis tipos de relación.}
    \label{fig:modelo-grafo}
\end{figure}

\begin{table}[H]
\centering
\caption{Elementos del grafo y su volumen al corte}
\label{tab:grafo}
\vspace{0.2cm}
\small
\begin{tabular}{@{}llp{5.6cm}r@{}}
\toprule
\textbf{Elemento} & \textbf{Tipo} & \textbf{Propiedades / origen} & \textbf{Cantidad} \\
\midrule
\texttt{:User}          & Nodo     & \texttt{login} (único)                            & 1,451,903 \\
\texttt{:Repo}          & Nodo     & \texttt{name} (único), \texttt{owner}             & 1,887,722 \\
\texttt{:Org}           & Nodo     & \texttt{login} (único)                            &   129,068 \\
\midrule
\texttt{PUSHED}         & Relación & \texttt{PushEvent}                                & 1,531,945 \\
\texttt{OPENED\_PR}     & Relación & \texttt{PullRequestEvent}                         &   654,560 \\
\texttt{STARRED}        & Relación & \texttt{WatchEvent}                               &   216,598 \\
\texttt{OPENED\_ISSUE}  & Relación & \texttt{IssuesEvent}                              &   175,188 \\
\texttt{FORKED}         & Relación & \texttt{ForkEvent}                                &    35,979 \\
\texttt{BELONGS\_TO}    & Relación & campo \texttt{org} del evento                     &   288,987 \\
\midrule
\multicolumn{3}{@{}l}{\textbf{Total de nodos}}      & \textbf{3,468,693} \\
\multicolumn{3}{@{}l}{\textbf{Total de relaciones}} & \textbf{2,903,257} \\
\bottomrule
\end{tabular}
\end{table}

Tres decisiones definen el modelo:

\begin{itemize}[leftmargin=*]
    \item \textbf{La relación agrega, no enumera.} Entre un usuario y un repositorio existe a lo sumo una relación de cada tipo, con la propiedad \texttt{count} acumulando cuántas veces ocurrió y \texttt{first\_seen}/\texttt{last\_seen} delimitando el periodo. Los 12.1 millones de \textit{pushes} con relación en el grafo se condensan así en 1.5 millones de aristas \texttt{PUSHED}. Modelar una arista por evento multiplicaría el tamaño del grafo sin aportar nada a las preguntas de conectividad.
    \item \textbf{El tipo de acción es el tipo de la relación.} Usar \texttt{PUSHED}, \texttt{STARRED}, etc.\ como tipos, en lugar de una relación genérica con una propiedad \texttt{tipo}, permite filtrar por tipo durante el recorrido sin leer propiedades.
    \item \textbf{Solo cinco tipos de evento generan relaciones.} Son los que expresan un vínculo directo de un usuario con un repositorio. Comentarios, revisiones y borrados no crean aristas.
\end{itemize}

\subsection{Restricciones}

Las tres restricciones de unicidad garantizan que no existan nodos duplicados y crean, a la vez, el índice que usa \texttt{MERGE} para localizar cada nodo:

\begin{lstlisting}[style=plano]
CREATE CONSTRAINT user_login IF NOT EXISTS FOR (u:User) REQUIRE u.login IS UNIQUE;
CREATE CONSTRAINT repo_name  IF NOT EXISTS FOR (r:Repo) REQUIRE r.name  IS UNIQUE;
CREATE CONSTRAINT org_login  IF NOT EXISTS FOR (o:Org)  REQUIRE o.login IS UNIQUE;
\end{lstlisting}

\subsection{Construcción del Grafo}

La etapa 5 (\repo{src/load_neo4j.py}{src/load\_neo4j.py}) no lee el archivo de origen ni MongoDB: lee la partición de la hora en la tabla \texttt{user\_repo\_\bk{}interactions\_\bk{}by\_hour} de Cassandra, agrupa las filas por tipo de relación y las envía en lotes de 1,000 con \texttt{UNWIND} y \texttt{MERGE}.

\bloque{Consulta de carga de relaciones (simplificada; el tipo de relación se sustituye por lote)}
\begin{lstlisting}[style=plano]
UNWIND $rows AS row
MERGE (u:User {login: row.user})
MERGE (r:Repo {name: row.repo})
  ON CREATE SET r.owner = split(row.repo, '/')[0]
MERGE (u)-[rel:PUSHED]->(r)
  ON CREATE SET rel.count = row.n,
                rel.first_seen = datetime(row.ts), rel.last_seen = datetime(row.ts)
  ON MATCH  SET rel.count = coalesce(rel.count, 0) + row.n,
                rel.last_seen = CASE WHEN datetime(row.ts) > rel.last_seen
                                     THEN datetime(row.ts) ELSE rel.last_seen END
\end{lstlisting}

\paragraph{Control de tamaño.} En cada hora solo entran al grafo los 10,000 repositorios y los 20,000 usuarios más activos de esa hora (parámetros \texttt{NEO4J\_TOP\_REPOS} y \texttt{NEO4J\_TOP\_USERS}). El grafo es, por tanto, una muestra sesgada hacia la actividad más intensa, lo que se declara en la metodología del \textit{dashboard}.

\paragraph{Idempotencia.} \texttt{MERGE} evita duplicar nodos y aristas, pero no evita sumar dos veces el \texttt{count}. Por eso la etapa consulta antes la columna \texttt{loaded\_neo4j} de \texttt{pipeline\_runs}: si la hora ya fue cargada en el grafo, se omite.

\paragraph{Reconstrucción.} Como el grafo se deriva íntegramente de Cassandra, puede regenerarse desde cero con \repo{scripts/rebuild_neo4j.py}{scripts/rebuild\_neo4j.py} sin volver a descargar ni a procesar los archivos de origen. Es la consecuencia práctica de la arquitectura por capas.

\subsection{Consultas de Grafo}

Las dos consultas siguientes ilustran por qué este dato reside en Neo4j. La primera es la base del KPI~5, que le añade un filtro por fecha sobre \texttt{first\_seen} (Sección~\ref{sec:kpis}); la segunda encuentra a los usuarios que comparten repositorios con uno dado, un recorrido de dos saltos.

\begin{lstlisting}[style=plano]
// Repositorios con más usuarios distintos conectados (grado del nodo)
MATCH (u:User)-[]->(r:Repo)
RETURN r.name AS repo, count(DISTINCT u) AS usuarios
ORDER BY usuarios DESC LIMIT 5;

// Usuarios que colaboran en los mismos repositorios que un usuario dado
MATCH (a:User {login: $login})-[]->(r:Repo)<-[]-(b:User)
WHERE a <> b
RETURN b.login AS colega, count(DISTINCT r) AS repos_en_comun
ORDER BY repos_en_comun DESC LIMIT 5;
\end{lstlisting}

\begin{table}[H]
\centering
\caption{Resultado de la primera consulta sobre el grafo acumulado}
\label{tab:top-repos}
\vspace{0.2cm}
\small
\begin{tabular}{@{}lr@{}}
\toprule
\textbf{Repositorio} & \textbf{Usuarios distintos conectados} \\
\midrule
\texttt{morluto/rea}                     & 6,843 \\
\texttt{storytold/photocraft}            & 1,483 \\
\texttt{Ebony-Vinyl/dsh-our-free-model}  & 1,148 \\
\texttt{boykopovar/AnyPS5}               & 1,025 \\
\texttt{debpalash/VoiceStudio}           &   842 \\
\bottomrule
\end{tabular}
\end{table}

El primer repositorio cuadruplica al segundo. El grafo permite ir más allá del número, porque conserva el tipo de cada relación (Tabla~\ref{tab:morluto}).

\begin{table}[H]
\centering
\caption{Relaciones entrantes de \texttt{morluto/rea} por tipo}
\label{tab:morluto}
\vspace{0.2cm}
\small
\begin{tabular}{@{}lrr@{}}
\toprule
\textbf{Tipo de relación} & \textbf{Relaciones} & \textbf{Interacciones (\texttt{count})} \\
\midrule
\texttt{STARRED}       & 5,628 & 5,630 \\
\texttt{FORKED}        & 1,156 & 1,156 \\
\texttt{OPENED\_PR}    &    58 &   143 \\
\texttt{OPENED\_ISSUE} &    26 &   173 \\
\texttt{PUSHED}        &     7 &   410 \\
\bottomrule
\end{tabular}
\end{table}

El 82\,\% de los usuarios conectados (5,608) tiene como única relación una estrella, y la actividad de código se reduce a 7 cuentas que hicieron \textit{push}. Un grado alto no implica por sí solo una comunidad de colaboradores: aquí describe popularidad, no colaboración, y esa distinción solo es visible al examinar el tipo de las relaciones. El vecindario se explora de forma interactiva en el \textit{dashboard} (Figura~\ref{fig:dash-neo4j-vecindario}).
